% Slides — Capítulo 12: Massas de Ar e Frentes
\documentclass[aspectratio=169,11pt]{beamer}
\usepackage{../comum/temameteo}
\graphicspath{{../figuras/cap12/}}

\title[Cap. 12 --- Massas de Ar e Frentes]{Capítulo 12 --- Massas de Ar e
Frentes}
\subtitle{Meteorologia Prática}
\author{}
\date{}

\begin{document}

\begin{frame}
  \titlepage
  \vfill
  \atribuicao
\end{frame}

\begin{frame}{Roteiro}
  \tableofcontents
\end{frame}

% ---------------------------------------------------------------
\section{Massas de ar}

\begin{frame}{O palco: ciclone, massas e frentes}
  \begin{columns}
    \column{0.4\textwidth}
    \begin{itemize}
      \item Massa de ar $=$ ar que estacionou e ganhou a cara da
            superfície; etiqueta: $(\theta_w, r)$ (Cap.~4).
      \item Frente $=$ zona de batalha entre duas massas.
      \item O ciclone (Cap.~13) é a onda que as põe em movimento.
    \end{itemize}
    \column{0.6\textwidth}
    \figlivroslide{fig_12_1.jpg}{0.54\textheight}{Fig.~12.1 --- ciclone
      idealizado com massas cP e mT e suas frentes}
  \end{columns}
\end{frame}

\begin{frame}{Onde nascem: os berçários do planeta}
  \begin{columns}
    \column{0.4\textwidth}
    \begin{itemize}
      \item Anticiclones estacionários: céu limpo, ar raso e parado —
            o chão carimba o ar (Caps.~10--11).
      \item Classificação de Bergeron: cP, cT, mP, mT ($+$ A, E).
      \item A friagem paulista: cP/mP andina modificada em trânsito.
    \end{itemize}
    \column{0.6\textwidth}
    \figlivroslide{fig_12_4.jpg}{0.52\textheight}{Fig.~12.4 --- regiões
      de gênese e tipos de massas de ar}
  \end{columns}
\end{frame}

\begin{frame}{Gênese fria: a EDO no quadro e no notebook}
  \begin{columns}
    \column{0.44\textwidth}
    \eqdestaque{$\rho C_p H\,\dfrac{dT}{dt} =
      -\varepsilon\sigma(T^4 - T_{sfc}^4)$}
    \vspace{0.4em}
    \begin{itemize}
      \item $\tau = \rho C_pH/4\varepsilon\sigma T^3 \sim$ dias:
            rápido no início, satura no fim (T2).
      \item Massa cP ``pronta'' em 3 dias a 2 semanas (Caso~1; N1).
      \item Gênese quente: subsidência $+$ mistura, de cima para
            baixo.
    \end{itemize}
    \column{0.56\textwidth}
    \figlivroslide{fig_12_5.jpg}{0.5\textheight}{Fig.~12.5 ---
      esfriamento da massa de ar em formação dia após dia}
  \end{columns}
\end{frame}

% ---------------------------------------------------------------
\section{Anatomia das frentes}

\begin{frame}{Frente fria em quatro camadas de carta}
  \begin{columns}
    \column{0.3\textwidth}
    \begin{itemize}
      \item A MESMA frente vista em isotermas, isóbaras (o cavado em
            ``V'', T4), ventos (o giro N$\to$SW no HS) e tempo
            presente.
      \item Treinar a sobreposição mental — é assim que se acha a
            frente na análise (Cap.~9).
    \end{itemize}
    \column{0.35\textwidth}
    \figlivroslide{fig_12_11a.jpg}{0.4\textheight}{Fig.~12.11 ---
      isotermas apertadas na zona frontal}
    \column{0.35\textwidth}
    \figlivroslide{fig_12_11b.jpg}{0.4\textheight}{Fig.~12.11 ---
      isóbaras dobrando no cavado frontal}
  \end{columns}
\end{frame}

\begin{frame}{Os cortes verticais clássicos}
  \begin{columns}
    \column{0.3\textwidth}
    \begin{itemize}
      \item Fria: cunha íngreme ($\sim$1/100), convecção na linha.
      \item Quente: rampa $\sim$1/200--1/300, estratiformes
            adiantados (Ci$\to$Cs$\to$As$\to$Ns).
      \item Inversão frontal marca a interface na sondagem (Cap.~5).
    \end{itemize}
    \column{0.35\textwidth}
    \figlivroslide{fig_12_14a.jpg}{0.4\textheight}{Fig.~12.14 --- frente
      fria: a cunha que empurra}
    \column{0.35\textwidth}
    \figlivroslide{fig_12_14b.jpg}{0.4\textheight}{Fig.~12.14 --- frente
      quente: a rampa que recua}
  \end{columns}
\end{frame}

\begin{frame}{Margules: a inclinação em equilíbrio}
  \begin{columns}
    \column{0.44\textwidth}
    \eqdestaque{$\tan\alpha = \dfrac{|f|\,\bar T\,\Delta U_g}
      {g\,\Delta T}$}
    \vspace{0.4em}
    \begin{itemize}
      \item Interface só estaciona com salto de vento ao longo da
            frente (T1).
      \item $\Delta T = 10$\,K, $\Delta U = 15$\,m/s: 1/230 — cunha
            quase deitada (Caso~2; N2).
      \item Sem rotação a cunha desaba: frentes são estruturas
            dinâmicas.
    \end{itemize}
    \column{0.56\textwidth}
    \figlivroslide{fig_12_17b.jpg}{0.5\textheight}{Fig.~12.17 --- o
      ajuste geostrófico da interface fria/quente}
  \end{columns}
\end{frame}

\begin{frame}{O jato mora em cima da frente}
  \begin{columns}
    \column{0.4\textwidth}
    \begin{itemize}
      \item Zona baroclínica concentrada $\Rightarrow$ vento térmico
            concentrado $\Rightarrow$ \alert{jato} na tropopausa
            (Cap.~11).
      \item O corte mostra as isotacas fechando no núcleo, bem sobre a
            zona frontal.
      \item CAT nas bordas (Cap.~5); streaks e ciclogênese (Cap.~13).
    \end{itemize}
    \column{0.6\textwidth}
    \figlivroslide{fig_12_19.jpg}{0.54\textheight}{Fig.~12.19 --- corte
      meridional: jato sobre a zona frontal}
  \end{columns}
\end{frame}

% ---------------------------------------------------------------
\section{Frontogênese}

\begin{frame}{Quem afia o gradiente}
  \begin{columns}
    \column{0.3\textwidth}
    \begin{itemize}
      \item Confluência, cisalhamento e inclinação: os três termos da
            função frontogenética.
      \item Deformação comprime isotermas no eixo de dilatação:
            $|\nabla T| \propto e^{at}$ (T3; Caso~3).
      \item Dobra a cada $\sim$19\,h — a frente nasce em 1--2 dias.
    \end{itemize}
    \column{0.35\textwidth}
    \figlivroslide{fig_12_22b.jpg}{0.42\textheight}{Fig.~12.22 ---
      confluência espremendo as isentropas}
    \column{0.35\textwidth}
    \figlivroslide{fig_12_26b.jpg}{0.42\textheight}{Fig.~12.26 ---
      deformação: isotermas dentro do ângulo crítico afiam}
  \end{columns}
\end{frame}

\begin{frame}{A circulação transversal: a quente sobe}
  \begin{columns}
    \column{0.4\textwidth}
    \begin{itemize}
      \item Frontogênese quebra o balanço térmico do vento; a
            atmosfera responde com uma circulação transversal
            \alert{termicamente direta}.
      \item Ar quente sobe sobre a cunha: a chuva frontal é filha
            desta célula.
      \item Mesma física dos quadrantes de jato (Cap.~13).
    \end{itemize}
    \column{0.6\textwidth}
    \figlivroslide{fig_12_28a.jpg}{0.5\textheight}{Fig.~12.28 ---
      circulação vertical transversal à zona frontal}
  \end{columns}
\end{frame}

% ---------------------------------------------------------------
\section{Oclusões, TROWAL e linha seca}

\begin{frame}{Oclusão: quem é mais frio fica embaixo}
  \begin{columns}
    \column{0.3\textwidth}
    \begin{itemize}
      \item Fria alcança a quente; setor quente erguido.
      \item Comparar $\theta_w$ dos dois lados frios (T5).
      \item TROWAL: a língua de ar quente erguida — faixa de nuvem e
            chuva enrolada no ciclone maduro (Cap.~13).
    \end{itemize}
    \column{0.35\textwidth}
    \figlivroslide{fig_12_30b.jpg}{0.44\textheight}{Fig.~12.30 ---
      oclusão de tipo frio em corte}
    \column{0.35\textwidth}
    \figlivroslide{fig_12_32.jpg}{0.44\textheight}{Fig.~12.32 --- o
      TROWAL nas isentropas}
  \end{columns}
\end{frame}

\begin{frame}{Frentes de altitude e a linha seca}
  \begin{columns}
    \column{0.3\textwidth}
    \begin{itemize}
      \item Dobra da tropopausa: ar estratosférico desce seco e rico
            em ozônio — escuro no canal do vapor (Cap.~8).
      \item Linha seca: fronteira de $T_d$ sem contraste de $T$ —
            gatilho clássico de severo (Cap.~14).
    \end{itemize}
    \column{0.35\textwidth}
    \figlivroslide{fig_12_35.jpg}{0.44\textheight}{Fig.~12.35 --- dobra
      da tropopausa junto ao jato}
    \column{0.35\textwidth}
    \figlivroslide{fig_12_37b.jpg}{0.44\textheight}{Fig.~12.37 --- a
      linha seca em corte vertical}
  \end{columns}
\end{frame}

% ---------------------------------------------------------------
\section{Síntese e laboratório}

\begin{frame}{O mapa do capítulo}
  \eqdestaque{anticiclone parado $\to$ massa $(\theta_w, r)$ $\to$
    contraste $\to$ frente:\quad
    $\tan\alpha = \tfrac{|f|\bar T\Delta U}{g\Delta T}$,\quad
    $|\nabla T|\propto e^{at}$}
  \vspace{0.6em}
  Massas são a memória térmica do chão; frentes, as costuras entre
  elas; a frontogênese fabrica a descontinuidade e a circulação
  transversal faz a chuva. O ciclone que anima tudo: Cap.~13.
\end{frame}

\begin{frame}{Laboratório numérico --- notebook do Cap.~12}
  \texttt{notebooks/cap12\_frentes.ipynb}
  \begin{enumerate}
    \item Gênese cP: a EDO radiativa do $T^4$.
    \item Margules: inclinação × contrastes.
    \item Frontogênese por deformação (mini-NWP com solução exata).
    \item Meteograma sintético da frente fria.
  \end{enumerate}
  \vspace{0.4em}
  \textbf{Exercícios}: T1--T5 e N1--N4 nas notas; células reservadas no
  notebook.
  \vfill
  \atribuicao
\end{frame}

\end{document}
